\section{Misure precise}

\subsection{Macchine di Turing precise}
Una macchina di Turing $M$ é \textbf{precisa} se esistono due funzioni $f$ e $g$ tali che:
\begin{itemize}
	\item Ogni computazione di $M$ termina \textit{esattamente} dopo $f(n)$ passi.
	\item Al termine di ogni computazione, tutti i nastri di $M$ sono lunghi \textit{esattamente} $g(n)$ (eccetto primo e ultimo nastro se la macchina é con input e output).
\end{itemize}
Formalmente: $\exists f, g: \mathbb{N} \rightarrow \mathbb{N}  \quad \forall x \in \Sigma^* : |x| = n$
\[
\begin{cases}
	\configurazionestart{x} \xrightarrow{{M}^t} \configurazione{q_H}{w}{u} \\
	\\
	t = f(n) \\
	|w_i u_i|= g(n) \quad \forall i \in \{2, \dots, k-1\} \\
	q_H \in \{h, "yes", "no"\}
\end{cases}
\]

\subsection{Funzioni di complessitá proprie}
Le funzioni di complessitá proprie sono funzioni $\phi$ sui naturali, non decrescenti e tali che una macchina di Turing precisa con input e output $M_\phi$ puó calcolarne il valore in unario, scrivendo sull'ultimo nastro tanti simboli $\sqcap$ ("quasi-blank") quanto vale $\phi(n)$. \\
Per qualsiasi input, il tempo impiegato deve essere sempre $f(n) = O(n + \phi(n))$ \\
Tutti i nastri devono usare sempre spazio $g(n) = O(\phi(n))$. \\

Formalmente: una funzione $\phi: \mathbb{N} \rightarrow \mathbb{N}$ é \textit{propria} sse:
\begin{itemize}
	\item $\phi$ non é decrescente: \quad $\forall n \in \mathbb{N} \quad \phi(n+1) \geq \phi(n)$
	\item Esiste una MdT precisa con input e output $M_\phi$ tale che:\\ $\forall x \in \Sigma^* : |x| = n$ \\
	\[
	\begin{cases}
		\configurazionestart{x} \xrightarrow{{M_\phi}^t} \configurazionepropria{x}{\phi(n)} \\
		\\
		t = f(n) = O(n + \phi(n)) \\
		z_i = g(n) = O(\phi(n)) \quad \forall i \in \{2, \dots, k-1\}
	\end{cases}
	\]
\end{itemize}

Le funzioni proprie sono le uniche funzioni che possono essere usate per misurare la complessitá di tempo e di spazio di un problema senza incorrere in problemi di misurazione.\\ 
Ad esempio, il Gap Theorem afferma che esistono funzioni che "rompono" l'analisi di complessitá.

\subsubsection*{Gap Theorem}
Esiste una funzione ricorsiva $f: \mathbb{N} \rightarrow \mathbb{N}$ tale che {$\TIME({f(n)}) = \TIME({2^{f(n)}})$}


\subsection{Teorema: potenza delle MdT con input e output}
Per ogni macchina di Turing con $k$ nastri $M$ che opera in tempo $f(n)$, esiste una macchina di Turing con input e output e $k+2$ nastri $M'$ che ha lo stesso output e opera in tempo $O(f(n))$.
\begin{proof}
M' copia l'input sul secondo nastro, simula M sui k nastri centrali, infine, se M ha un output, lo copia sull'ultimo nastro. \\
Il costo dell copia dell'input sul secondo nastro é $O(n)$. \\
Poiché sappiamo che $f(n)$ é limitato inferiormente da $n$, possiamo dire che il costo della copia dell'input é $O(f(n))$. \\
Il costo della simulazione di M é $O(f(n))$. \\
Poiché sappiamo che lo spazio é limitato superiormente da $f(n)$, il costo della copia dell'output sull'ultimo nastro é $O(f(n))$, . \\
Quindi il costo complessivo é $O(f(n))$.
\end{proof}
Di conseguenza, le macchine di Turing con input e output sono equipotenti a quelle normali.


\subsection{Teorema: soluzioni precise}
\subsubsection{Versione tempo}
Se una MdT $M$ \textit{decide} $L$ in tempo $f(n)$, ed $f$ é una funzione \textit{propria}, allora esiste una MdT \textit{precisa} $M'$ che decide $L$ in tempo $O(f(n))$.
\begin{proof}
Costruiamo $M'$ come una concatenazione di due macchine di Turing, tenendo i loro nastri separati. \\
La prima MdT é $M_f$, cioé la macchina con input e output che calcola la funzione propria $f$. \\
$M_f$ Scrive $\sqcap^{f(n)}$ sul nastro $k$  impiegando sempre tempo $O(n + f(n)) = O(f(n))$. \\
Quando termina, si passa a una macchina $M_t$, che simula $M$ sul proprio set di nastri $\{k+1, \dots, 2k\}$, ma a ogni passo sposta anche la testina sul nastro $k$ verso destra. \\
Sappiamo che $M$, e quindi anche $M_t$, termina entro $f(n)$ passi. Se dovesse terminare prima, basta aspettare che la testina sull'ultimo nastro finisca di scorrere tutti i $\sqcap$, impiegando esattamente $f(n)$ passi. \\
In totale, quindi, $M' = (M_f, M_t)$ impiega sempre $O(f(n)) + f(n) = O(f(n))$ passi.
\end{proof}

\subsubsection{Versione spazio}
Se una MdT $M$ \textit{decide} $L$ in spazio $g(n)$, e $g$ é una funzione \textit{propria}, allora esiste una MdT \textit{precisa} $M'$ che decide $L$ in spazio $O(g(n))$.
\begin{proof}
Costruiamo $M'$ come una concatenazione di due macchine di Turing, tenendo i loro nastri separati. \\
La prima MdT é $M_g$, cioé la macchina con input e output che calcola la funzione propria $g$. \\
$M_g$ Scrive $\sqcap^{g(n)}$ sul nastro $k$  usando spazio $O(g(n))$ \\
Quando termina, si passa a una macchina $M_s$, che copia tutti gli $\sqcap$ sul nastro $k$ su ogniuno dei suoi nastri $\{k+1, \dots, 2k\}$ e poi simula $M$ usando $\sqcap$ al posto di $\sqcup$. \\
$M$, e quindi $M_s$, non usano piú di $g(n)$ celle su ogni nastro. Se non sono stati tutti sovrascritti gli $\sqcap$ contano come celle occupate sul nastro, rendendo tutti i nastri lunghi $g(n)$.
\end{proof}

Questo teorema ci permette di usare macchine di Turing non precise per la misurazione della complessitá, a patto che le funzioni $f, g$ siano proprie.


\subsection{Teorema: speedup multinastro}
\label{teorema:speedup_multinastro}
Per ogni MdT $M$ a $k$ nastri che opera in tempo $f(n)$, esiste una MdT $M'$ a $1$ nastro che: 
\begin{itemize} 
	\item Opera in tempo $O(f(n)^2)$ 
	\item $\forall x \in \Sigma^* \quad M(x) = M'(x)$ 
\end{itemize}
\begin{proof}
Data $M = (Q, \Sigma, \delta, s)$. \\
Costruiamo $M' = (Q', \Sigma', \delta', s)$ che simula i $k$ nastri di $M$ usando un solo nastro. \\
Si separa il nastro di $M'$ in $k$ sezioni, introducendo nuovi simboli: 
\begin{itemize}
	\item $\blacktriangleright$: "pseudo-start", rappresenta gli $\triangleright$ dei $k$ nastri di $M$
	\item $\blacktriangleleft$: separatore dei $k$ nastri di $M$
	\item $\triangleleft$: rappresenta la fine della stringa di $M'$
	\item I simboli sottolineati $\underline{\Sigma} = \{\underline{\sigma} : \sigma \in \Sigma \}$ indicano cosa c'é sotto le testine di $M$.
\end{itemize}
{
Una configurazione di $M$:
\[
	\configurazione{q}{w}{u}
\] 
Corrisponde a una configurazione di $M'$:
\[
	q ({\triangleright}) (\blacktriangleright w_1 u_1 \blacktriangleleft \dots \blacktriangleright w_k u_k \blacktriangleleft \triangleleft )
\]
}
{
La lunghezza del nastro di $M'$ é al piú:
\begin{itemize}
	\item $k \cdot f(n)$ per i nastri di $M$ 
		\subitem ($M$ ha $k$ nastri lunghi al piú $f(n)$)
	\item $2k$ per i simboli $\blacktriangleright$ e $\blacktriangleleft$
	\item $2$ per i simboli $\triangleright$ e $\triangleleft$
\end{itemize}
Cioé $O(f(n))$. \\
}

La configurazione iniziale di $M'$ é: \[s(\triangleright)(x)\]
La preparazione del nastro si articola in:
\begin{enumerate}
	\item Traslare a destra di una posizione l'input $x$ e inserire $\blacktriangleright$ all'inizio.
		\begin{itemize}
			\item Per farlo si introducono $|\Sigma| + 1$ nuovi stati (vedere esempio \ref{example:shift})
			\item Il tempo impiegato é $O(n)$.
			\item Il nastro di $M'$ diventa: $\triangleright \blacktriangleright x$.
		\end{itemize} 
	\item Appendere la stringa che rappresenta gli altri nastri vuoti: $\blacktriangleleft ({\blacktriangleright \blacktriangleleft})^{k-1} \triangleleft$
		\begin{itemize}
			\item Si usano $2k$ nuovi stati, uno per ognuno dei due simboli, che tengono traccia del nastro corrente in modo da fermarsi a $k$. 
			\begin{itemize}
				\item Fa eccezione il primo nastro, per cui serve un solo stato che scriva $\blacktriangleleft$ e poi passi al nastro successivo. 
				\item Ma a questi si aggiunge un ultimo stato che scrive $\triangleleft$.
			\end{itemize}
			\item Il tempo impiegato é $O(k)$.
		\end{itemize}
\end{enumerate}
{
Dopo aver preparato il nastro, si puó iniziare la simulazione. \\
Ogni passo di $M$ richiede due step:
\begin{enumerate}
	\item Scansione del nastro per leggere i simboli sottolineati, cioé quelli sotto le testine di $M$, fino a $\triangleleft$.
		\begin{itemize}
			\item Servono stati che ricordano lo stato originale di $M$ e ogni simbolo letto finora. Per ogni stato $q \in Q$, aggiungo a $Q'$ tutti gli stati $q_{\sigma_1}; q_{\sigma_1, \sigma_2}; \dots; q_{\sigma_1, \dots, \sigma_k}$, per tutte le combinazioni di lunghezza minore o uguale a $k$ di simboli $\sigma_i \in \Sigma$. 
			\item In totale, servono $|Q| \cdot |\Sigma|^{1 + \dots + k} = |Q| \cdot |\Sigma|^{k \cdot (k+1)/2}$ nuovi stati.
			\item Il tempo	 impiegato é $O(f(n))$ (teorema \ref{teorema:spazio_limitato_superiormente_dal_tempo}).
		\end{itemize}
	\item Modifica del nastro per applicare la transizione a partire da uno stato $q_{\sigma_1,\dots,\sigma_k}$.
		\begin{itemize}
			\item Si scandisce il nastro all'indietro fino a un simbolo sottolineato, che viene modificato come avrebbe fatto $M$, ed eventualmente effettuando un passo a sinistra o a destra e sottolineando il nuovo simbolo.
			\item Se la testina si sposta a destra e si trova su $\blacktriangleleft$, allora bisogna effettuare uno shift di tutti i simboli a destra.
				\begin{enumerate}
					\item Si sovrascrive $\blacktriangleleft$ con $\overline{\blacktriangleleft}$
					\item Si va a destra fino a $\triangleleft$
					\item Tornando indietro, si effettua uno shift di tutti i simboli verso destra, fino a tornare a $\overline{\blacktriangleleft}$, che viene sovrascritto con $\underline{\sqcup}$ e riportato a destra come $\blacktriangleleft$.
					\item Il tempo dello shift é $O(f(n))$.
				\end{enumerate}
			\item Il tempo impiegato per la scansione é $O(f(n))$. Nel caso peggiore, bisogna effettuare uno shift per ogni passo, quindi il tempo totale é $O(f(n)^2)$.
		\end{itemize}
\end{enumerate}
}

Quando $M$ termina, $M'$ sposta il contenuto del nastro $k$ sul nastro $1$, cancellando tutto il resto. Questo spostamento impiega tempo $O(f(n))$. Poi termina.

Il termine dominante del tempo é quello della modifica del nastro durante la simulazione: $O(f(n)^2)$. \\
\end{proof}

Questo teorema ci permette di usare quanti nastri vogliamo per la misurazione della complessitá, la classe di complessitá del problema non cambia perché lo speedup é al piú quadratico. \\




\subsection{Teorema: Linear Speedup}
\subsubsection{Versione tempo}
Se $L \in \TIME(f(n))$ allora $\forall \varepsilon > 0 \quad L \in \TIME(f'(n))$ \\
Dove $f'(n) = \varepsilon f(n) + n + 2$
\begin{proof}
Simuliamo $M=(Q, \Sigma, \delta, s)$ con una macchina $M'=(Q', \Sigma', \delta', s')$. \\
L'idea é simulare $m$ passi di $M$ con $6$ passi di $M'$, scegliendo $m$ opportunamente. \\
Per farlo, occorre utilizzare un numero di stati e di simboli dell'alfabeto esponenziale rispetto a quelli della macchina originale. \\ 
$\Sigma' = \Sigma \cup \Sigma^m$: contiene tutti i simboli dell'alfabeto originale e tutte le $m$-uple di simboli su quell'alfabeto.
Cioé consideriamo $(\sigma_1, \ldots, \sigma_m)$ come un \textit{unico} simbolo. \\

La macchina comincia con una scansione del primo nastro per comprimere l'input: ogni sequenza di $m$ simboli di $x$ viene scritta con un solo simbolo di $\Sigma^m$ sul secondo nastro.\\
Opzionalmente, si puó cancellare $x$ durante la scansione per avere un nastro in piú a disposizione. \\
Notare che se $M$ ha un solo nastro, allora $M'$ ha bisogno di due nastri. In tutti gli altri casi possono avere lo stesso numero di nastri. \\
\begin{itemize}
    \item dallo stato iniziale $s'$ usiamo stati da ${s'}_{\sigma_1}$ a ${s'}_{\sigma_1, \ldots, \sigma_{m-1}}$ per ricordare ogni simbolo incontrato (inclusi i $\sqcup$ se $n$ non é divisibile per $m$).
    \item in uno degli stati che ricordano $m-1$ simboli, leggiamo quello sotto la testina e scriviamo sul secondo nastro $(\sigma_1 \ldots, \sigma_m)$.
    \item Tempo: $m \cdot \ceil{n/m} + 2$  passi sul primo nastro, piú $\ceil{n/m}$ sul secondo.
\end{itemize}

Ora la simulazione comincia usando la stringa del secondo nastro come input. \\
Verranno utilizzati \textit{solo} simboli di $\Sigma^m$ per simulare $m$ passi alla volta. \\

\textbf{Se la testina del nastro $i$ di $M$ si trova sulla cella $l_i$, allora la testina del nastro $i$ di $M'$ si trova sulla cella $\ceil{l_i/m}$.} \\
Di conseguenza, per conoscere la posizione della testina di $M$ sul nastro $i$, basta ricordare solo in quale posizione sta tra le $m$ del simbolo corrente di $M'$. \\
Il numero di stati richiesti per ricordare lo stato di $M$ e la posizione della testina é $|Q| \cdot m^k$. \\

Se $M$ é nello stato $q$ e la testina del suo nastro $i$ si trova sulla cella $l_i$, allora chiamiamo ${l'}_i = (l_i \mod m)$ e diciamo che $M'$ é nello stato $q_{l'_1, \ldots, l'_k}$. .

Simulando $m$ passi di $M$, il cursore di un nastro di $M'$ potrebbe finire in una delle due caselle adiacenti a quella attuale, e non oltre (perché lo spazio é limitato superiormente dal tempo).\\
Per questo, si leggono i simboli a sinistra e a destra del cursore (1 passo a sinistra, 2 a destra e 1 a sinistra tornando alla posizione originaria: 4 in totale). \\
$\delta'$ é definita in modo da applicare direttamente le modifiche al blocco centrale e, eventualmente, uno dei due adiacenti. Questo richiede al massimo 2 passi. \\

In totale, quindi, ogni simulazione di $m$ passi di $M$ richiede $6$ passi di $M'$. \\
$M'$ impiega $6 \cdot \ceil{f(n)/m} + n + 2$ passi per simulare $M$. \\
Imponendo $m = \ceil{6/\varepsilon}$ otteniamo $\varepsilon f(n) + n + 2$.
\end{proof}

Questo teorema ci permette di ignorare costanti moltiplicative nella misurazione della complessitá di tempo.

\subsubsection{Versione spazio}
Se $L \in \SPACE(f(n))$ allora $\forall \varepsilon > 0 \quad L \in \SPACE(f'(n))$ \\
Dove $f'(n) = \varepsilon f(n) + n + 2$

Questo teorema ci permette di ignorare costanti moltiplicative nella misurazione della complessitá di spazio.
